% =============================================================================
% Kapitel 2: Notwendigkeit und Anwendungsfelder von Off-Grid Schutzsensorik
% =============================================================================

\chapter{Einleitung und Motivation}
\label{chap:motivation}

Dieses Kapitel legt die anwendungsseitige Grundlage der Arbeit. Zunächst werden
repräsentative Einsatzszenarien analysiert, um die spezifischen Anforderungen an
Schutzsensorik in Gefahrgebieten zu definieren. Darauf aufbauend wird dargelegt,
warum konventionelle Kommunikationstechnologien in diesen Umgebungen strukturell
versagen. Abschließend erfolgt ein Vergleich von \acrshort{lpwan}"=Technologien, um die
Wahl von LoRaWAN für das in dieser Arbeit entwickelte System zu begründen.


% -----------------------------------------------------------------------------
\section{Mögliche Einsatzszenarien}
\label{sec:einsatzszenarien}
% -----------------------------------------------------------------------------

\subsection{Industrie und Bergbau}
\label{subsec:uc-industrie}

Der Untertagebau gehört zu den gefährlichsten Arbeitsbereichen weltweit. Gesteinsschichten
verursachen eine extreme Signaldämpfung, die herkömmliche Funksysteme bereits nach
wenigen Metern unbrauchbar macht, während gleichzeitig die kontinuierliche Überwachung
toxischer und explosiver Gase, insbesondere Methan und Kohlenmonoxid,
überlebenswichtig ist. Methan akkumuliert in schlecht belüfteten Stollen
und kann bei Erreichen der unteren Explosionsgrenze von \SI{5}{\percent} Volumenanteil
in Luft zu verheerenden Explosionen führen~\cite{rumbleCRCHandbookChemistry2021}. Ein Ausfall der
Gasüberwachung in explosionsgefährdeten Zonen nach ATEX"=Richtlinie 2014/34/EU
kann damit katastrophale Folgen haben.

Auch in chemischen Anlagen und Raffinerien ist die zuverlässige Detektion von
Schadstoffleckagen eine sicherheitskritische Aufgabe. Weitläufige Industrieareale
mit Hindernissen, Metallstrukturen und variierenden Umgebungsbedingungen
erschweren den Aufbau einer zuverlässigen Kommunikationsinfrastruktur zusätzlich.

\subsection{Katastrophenschutz und Rettungswesen}
\label{subsec:uc-katastrophenschutz}

Einsatzkräfte des Katastrophenschutzes und der Feuerwehr operieren per Definition
in Umgebungen, in denen die Infrastruktur entweder nicht vorhanden oder durch
das Ereignis selbst zerstört worden ist. Bei Erdbeben werden Basisstationen und
Stromversorgung häufig vom selben Ereignis getroffen wie die Einsatzkräfte; bei
Waldbränden befinden sich Helfer in weitläufigen, abgelegenen Gebieten ohne
Mobilfunkabdeckung~\cite{sciulloDesignPerformanceEvaluation2020}.

In diesen Szenarien ist ein rasch deployierbares, autarkes Kommunikationsnetz
erforderlich, das ohne externe Stromversorgung oder Internetanbindung funktioniert.
Neben der reinen Datenübertragung von Vitaldaten und Schadstoffmesswerten ist
die Personenverfolgung eine wesentliche Funktion: die Einsatzleitung muss nämlich jederzeit wissen, wo sich einzelne 
Truppen befinden, um im Gefahrenfall koordiniert reagieren zu können. 

\subsection{Umweltmonitoring in Extremgebieten}
\label{subsec:uc-umwelt}

Bei der Überwachung von Vulkanen, Gletschern oder Hochwasser"=Pegelständen
in infrastrukturschwachen Regionen müssen Sensorknoten über Monate oder Jahre
hinweg energieautark arbeiten. Physischer Zugang für Wartungsarbeiten ist oft
lebensgefährlich oder logistisch unmöglich, denn ein Sensorknoten an einem Vulkanhang etwa
kann nicht für einen regulären Batteriewechsel aufgesucht werden.

Dies setzt extrem niedrige mittlere Stromverbräuche voraus, die nur durch
konsequente Nutzung von Tiefschlafmodi und ereignis"= oder zeitgesteuerte
Übertragungsstrategien erreichbar sind. Gleichzeitig darf die Reichweite nicht
durch den Energiesparbetrieb so weit eingeschränkt werden, dass die Datenübertragung
an die (oft ebenfalls entfernte) Basisstation nicht mehr zuverlässig möglich ist.

\subsection{Operative Anforderungen an die Schutzsensorik}
\label{subsec:anforderungen-zusammenfassung}

Aus den drei analysierten Szenarien lassen sich die zentralen technischen
Anforderungen ableiten, die ein geeignetes Kommunikationssystem für
Off"=Grid"=Schutzsensorik erfüllen muss: 

\begin{table}[htbp]
    \centering
    \caption{Vergleich der Einsatzszenarien hinsichtlich operativer Anforderungen}
    \label{tab:usecases}
    \begin{tabularx}{\textwidth}{l X X X}
        \toprule
        & \textbf{Industrie /} & \textbf{Katastrophen-}
        & \textbf{Umwelt-} \\
        & \textbf{Bergbau} & \textbf{schutz} & \textbf{monitoring} \\
        \midrule
        Infrastruktur
            & Keine (Untertagebau)
            & Zerstört / keine
            & Keine \\
        \addlinespace
        Reichweite
            & \SIrange{100}{500}{\m}
            & \SIrange{500}{2000}{\m}
            & \SIrange{2}{15}{\km} \\
        \addlinespace
        Batterielaufzeit
            & 8--24\,h
            & 4--48\,h
            & Monate--Jahre \\
        \addlinespace
        Latenztoleranz
            & Niedrig ($<$\,\SI{5}{\s})
            & Mittel ($<$\,\SI{30}{\s})
            & Hoch ($<$\,\SI{30}{\min}) \\
        \addlinespace
        Mobilität der Knoten
            & Hoch
            & Hoch
            & Keine \\
        \bottomrule
    \end{tabularx}
\end{table}

Die drei übergeordneten Anforderungen, die alle Szenarien gemeinsam haben, sind damit: 
\begin{itemize}
    \item \textbf{Infrastrukturunabhängigkeit}: das Netz muss ohne externe
    Basisstationen oder Internetanbindung vollständig funktionsfähig sein.
    \item \textbf{Hohe Reichweite bei starker Dämpfung}: das Signal muss Gestein,
    Beton und Vegetation durchdringen.
    \item \textbf{Maximale Energieeffizienz}: der Stromverbrauch der Endgeräte
    muss einen Akkubetrieb über den gesamten Einsatzzeitraum ermöglichen.
\end{itemize}

% -----------------------------------------------------------------------------
\section{Grenzen konventioneller Kommunikationsinfrastrukturen}
\label{sec:komm-versagen}
% -----------------------------------------------------------------------------

Gängige drahtlose Kommunikationstechnologien wie \acrshort{wlan}, Bluetooth und
Mobilfunk der vierten und fünften Generation sind für den Einsatz in erschlossenen,
infrastrukturell versorgten Umgebungen optimiert. Gegenüber den in
Abschnitt~\ref{sec:einsatzszenarien} definierten Anforderungen weisen sie
strukturelle Defizite auf, die einen zuverlässigen Einsatz in Gefahrgebieten
ausschließen.

\begin{itemize}
    \item \textbf{Reichweite:} WLAN nach IEEE\,802.11 erreicht im
    Innenbereich typischerweise 30 bis \SI{100}{\m} und erfordert die Installation
    von Access"=Points in regelmäßigen Abständen -- eine Voraussetzung, die im
    Untertagebau oder in Katastrophengebieten nicht erfüllbar ist. Bluetooth ist
    mit einer Reichweite von bis zu \SI{100}{\m} (Klasse\,1) für weiträumige
    Gefahrgebiete ebenfalls ungeeignet.

    \item \textbf{Infrastrukturabhängigkeit:} Mobilfunknetze bieten zwar eine
    großflächige Abdeckung, weisen jedoch einen systemimmanenten
    \textit{Single"=Point"=of"=Failure} auf: Basisstationen werden häufig durch
    dasselbe Ereignis zerstört oder überlastet, das den Notfall verursacht, wie zum Beispiel
    Erdbeben, Überschwemmungen oder Brände~\cite{sciulloDesignPerformanceEvaluation2020}. Private % TODO: maybe remove this citation
    \acrshort{lte}"= oder 5G"=Campusnetze könnten diese Abhängigkeit prinzipiell
    aufheben, erfordern jedoch lizenziertes Spektrum, einen dedizierten
    Core"=Network"=Betrieb und eine erhebliche Infrastrukturinvestition, die
    für mobile Einsatzszenarien nicht praktikabel ist.

    \item \textbf{Energieverbrauch:} Die für hohe Datenraten optimierten
    Protokolle, insbesondere LTE und 5G, verursachen einen hohen
    Strombedarf, der den Langzeitbetrieb kleiner, batteriebetriebener Sensorknoten
    unmöglich macht. Ein LTE"=Modul im Sendebetrieb verbraucht typischerweise
    mehrere hundert Milliampere, was bei einer \SI{1000}{mAh}"=Zelle eine Betriebsdauer von wenigen Stunden ergäbe.
\end{itemize}

% TODO: fact check, especially when it comes to LTE
\begin{table}[htbp]
    \centering
    \caption{Eignung gängiger Funktechnologien für Off-Grid-Gefahrgebiete}
    \label{tab:tech_vergleich_kurz}
    \begin{tabularx}{\textwidth}{l X X X X}
        \toprule
        \textbf{Technologie} & \textbf{Reichweite} & \textbf{Infra"-struktur"-frei}
        & \textbf{Energie"-effizienz} & \textbf{Geeignet} \\
        \midrule
        WLAN (802.11)
            & \SIrange{30}{100}{\m}
            & Nein  & Mittel & Nein \\
        \acrshort{ble} (Bluetooth)
            & bis \SI{100}{\m}
            & Ja    & Hoch   & Nein \\
        LTE / 5G
            & \SIrange{1}{10}{\km}
            & Nein  & Gering & Nein \\
        \bottomrule
    \end{tabularx}
\end{table}

Aus Tabelle~\ref{tab:tech_vergleich_kurz} ist ersichtlich, dass keine konventionelle Kommunikationstechnologie
alle drei genannten Kernanforderungen gleichzeitig erfüllen können. Für den Einsatz in Gefahrgebieten bedarf es also 
einer Alternative. % TODO: maybe polish this transition sentence

% -----------------------------------------------------------------------------
\section{LPWAN als Lösungsansatz für Off-Grid-Szenarien}
\label{sec:lpwan-solution}
% -----------------------------------------------------------------------------
% TODO: write here

\subsection{Vergleich prominenter LPWAN-Technologien}
\label{subsec:lpwan-vergleich}

Im \acrshort{lpwan}"=Segment existieren drei etablierte Standards: Sigfox, \acrshort{nbiot}
und LoRaWAN. Obwohl alle drei auf niedrige Datenraten und langen Akkubetrieb
ausgelegt sind, unterscheiden sie sich fundamental in ihrer Netzarchitektur
und damit in ihrer Eignung für Off"=Grid"=Szenarien~\cite{mekki2019comparative}.

\textbf{Sigfox} betreibt ein proprietäres, öffentliches Netzwerk, das auf
der Infrastruktur des Unternehmens Sigfox S.A. basiert. Endgeräte können
ausschließlich über dieses öffentliche Netz kommunizieren; der Aufbau eines
privaten, autarken Netzes ist technisch nicht vorgesehen. Für
Off"=Grid"=Szenarien ohne Internetanbindung scheidet Sigfox damit aus.

\textbf{\acrshort{nbiot}} ist ein von 3GPP standardisiertes \acrshort{lpwan}"=Protokoll,
das im lizenzierten Mobilfunkspektrum (LTE"=Band) operiert und von
Mobilfunkbetreibern als Managed Service angeboten wird. Obwohl \acrshort{nbiot}
exzellente Gebäudedurchdringung und garantierte Dienstgüte (\textit{Quality
of Service}) bietet, erfordert es zwingend eine aktive SIM"=Karte und eine
Verbindung zum Kernnetz des Mobilfunkbetreibers~\cite{mekki2019comparative}.
In Gebieten ohne Mobilfunkabdeckung -- dem Kernszenario dieser Arbeit -- ist
\acrshort{nbiot} daher nicht einsetzbar.

\textbf{LoRaWAN} operiert im lizenzfreien \acrshort{ism}"=Band und ermöglicht als
einzige der drei Technologien den vollständig autarken Betrieb eines privaten
Netzwerks ohne Internetanbindung und ohne Abhängigkeit von einem
Netzwerkbetreiber. Gateway, Network"=Server und Application"=Server können
lokal -- etwa auf einem batteriebetriebenen Einplatinenrechner -- betrieben
werden. Tabelle~\ref{tab:lpwan_vergleich} fasst den Vergleich zusammen.

\begin{table}[htbp]
    \centering
    \caption{Technischer Vergleich führender LPWAN-Standards}
    \label{tab:lpwan_vergleich}
    \begin{tabularx}{\textwidth}{l X X X}
        \toprule
        \textbf{Merkmal} & \textbf{Sigfox} & \textbf{NB-IoT} & \textbf{LoRaWAN} \\
        \midrule
        Spektrum
            & Unlizenziert
            & Lizenziert (LTE)
            & Unlizenziert \\
        \addlinespace
        Betriebsart
            & Nur öffentlich
            & Nur öffentlich
            & \textbf{Öffentlich oder privat} \\
        \addlinespace
        Internetanbindung erforderlich
            & Ja
            & Ja
            & \textbf{Nein} \\
        \addlinespace
        Autarker Netzbetrieb möglich
            & Nein
            & Nein
            & \textbf{Ja} \\
        \addlinespace
        Typische Reichweite
            & \SIrange{3}{10}{\km}
            & \SIrange{1}{10}{\km}
            & \SIrange{2}{15}{\km} \\
        \addlinespace
        Off-Grid-Eignung
            & Nein
            & Nein
            & \textbf{Hervorragend} \\
        \bottomrule
    \end{tabularx}
\end{table}

\subsection{Begründung der Technologiewahl: LoRaWAN}
\label{subsec:technologiewahl}

Für das in dieser Arbeit entwickelte System wurde LoRaWAN gewählt, da es als
einzige der betrachteten Technologien den Aufbau eines vollständig privaten
Netzwerks ohne Internetverbindung ermöglicht. Als Network"=Server kommt
\textbf{ChirpStack} zum Einsatz~\cite{chirpstack_docs}, da dieser quelloffen,
auf ressourcenarmer Hardware (z.\,B. Raspberry Pi) betreibbar und vollständig
offline"=fähig ist -- im Gegensatz zu The Things Stack, dessen kostenlose
Community"=Edition eine permanente Verbindung zu den TTN"=Cloud"=Servern erfordert.

Darüber hinaus ermöglicht die Integration des LoRa"=Transceivers im
\textbf{STM32WL}"=System"=on"=Chip eine hochintegrierte und energieeffiziente
Hardwarebasis, bei der kein separates Funkmodul erforderlich ist. Die technischen
Grundlagen von LoRa und LoRaWAN werden im folgenden Kapitel~\ref{chap:grundlagen}
ausführlich behandelt.

\newacronym{wlan}{WLAN}{Wireless Local Area Network}
\newacronym{lpwan}{LPWAN}{Low Power Wide Area Network}
\newacronym{lte}{LTE}{Long Term Evolution}
\newacronym{ble}{BLE}{Bluetooth Low Energy}
\newacronym{nbiot}{NB-IoT}{Narrow Band Internet of Things}